
% 附录源 PDF 第 22—24 页；A.4.2。
\subsection{基于有效影响曲线各分量条件期望的目标参数一阶展开}

有
\[
D^N_{g^*}(P)(O^N)=\frac1N\sum_{i=1}^N\bar D_{g^*}(P)(\bar O(i)).
\]
将上面的恒等式与解有效影响曲线方程
$D^N_{g^*}(P_N^*)(O^N)=0$ 的 TMLE 结合，可得
\[
\begin{aligned}
\Psi_{g^*}(P_N^*)-\Psi_{g^*}(P_0)
={}&\frac1N\sum_{i=1}^N\{\bar D_{g^*}(P_N^*)(\bar O(i))-P_0^N\bar D_{g^*})\}\\
&+R_{21,g^*}(P_N^*,P_0).
\end{aligned}
\]
利用鞅理论，可以在合理假设下证明
\[
\frac1{\sqrt N}\sum_{i=1}^N
 \{\bar D_{g^*}(P_N^*)(\bar O(i))-P_{0,o(i)}\bar D_{g^*}(P_N^*)\}
\]
收敛到正态极限分布。因此，还剩下余项
\[
R_{22,g^*,N}(P_N^*,P_0)
 =\frac1N\sum_{i=1}^N(P_{0,o(i)}-P_0^N)\bar D_{g^*}(P_N^*),
\]
我们希望证明它是二阶余项，从而可以合理地假设其为 $o_P(N^{-1/2})$。有
\[
\bar D_{g^*}(P)=\sum_{x\in\{y,a,w\}}\sum_{s=1}^{\tau}\bar D_{g^*,x(s)}(P),
\]
其中 $\bar D_{g^*,x(s)}$ 是有效影响曲线中对应于给定 $C_x(s)$ 时 $X(s)$ 条件分布得分的分量。因此，
\[
R_{22,g^*,N}(P_N^*,P_0)
 =\sum_{x\in\{y,a,w\}}\sum_{s=1}^{\tau}R_{22,x(s),g^*,N}(P_N^*,P_0),
\]
其中
\[
R_{22,x(s),g^*,N}(\theta,\theta_0)
 =\frac1N\sum_{i=1}^N(P_{0,o(i)}-P_0^N)\bar D_{g^*,x(s)}(P).
\]
所以，只需研究这个与 $x(s)$ 对应的项。定义
\[
\begin{aligned}
&f_{x(s),g^*}(\theta,\theta_0)(C_x(i))\\
&\equiv\frac{h^*_{c_x(s)}(C_x(i))}{\bar h_{c_x}(C_x(i))}
 \int_x\Phi_{x(s)}(x,C_x(i))
 \,\mathrm d(P^*_{x(s)}-P^*_{0,x(s)})(x\mid C_x(i)),
\end{aligned}
\]
注意，这是 $P$ 与 $P_0$ 之间的一阶差异。另外，
\[
\begin{aligned}
&E_{P_0}f_{x(s),g^*}(\theta,\theta_0)(C_x(i))\\
&=\int_x\int_c\frac{\bar h_{c_x,0}(c)}{\bar h_{c_x}(c)}
 \Phi_{x(s)}(x,c)\,\mathrm d(P^*_{x(s)}-P^*_{0,x(s)})(x\mid c)
 h^*_{c_x(s)}(c)\,\mathrm d\mu(c).
\end{aligned}
\]
有
\[
\begin{aligned}
R_{22,x(s),g^*,N}(\theta,\theta_0)
=\frac1N\sum_{i=1}^N\Bigl\{
&f_{x(s),g^*}(\theta,\theta_0)(C_x(i))\\
&-E_0f_{x(s),g^*}(\theta,\theta_0)(C_x(i))\Bigr\}.
\end{aligned}
\]
因此，$R_{22,x(s),g^*,N}(\theta,\theta_0)$ 是 $N$ 个均值为零的随机变量的平均，并且每个随机变量都在 $P\to P_0$ 时趋于零。它涉及双重差异，所以是二阶余项。结合前面的定理 A.5，得到下列定理。

\begin{bookresult}{定理}{A.4}
令 $\Phi_{x(s)}(x,c_x)=E(Y_{g^*}(\tau)\mid X(s)=x,C_x(s)=c_x)$。回顾有效影响曲线的表达
$D^N_{g^*}(P)=N^{-1}\sum_{i=1}^N\bar D_{g^*}(P)(\bar O(i))$。
有
\[
\begin{aligned}
-\frac1N\sum_{i=1}^NP_{0,o(i)}\bar D_{g^*}(P)
={}&\Psi_{g^*}(P)-\Psi_{g^*}(P_0)\\
&+R_{21,g^*}(\theta,\theta_0)+R_{22,g^*,N}(\theta,\theta_0),
\end{aligned}
\]
其中 $R_{21,g^*}(\theta,\theta_0)$ 在定理 A.5 中定义，且
\[
R_{22,g^*,N}(\theta,\theta_0)
 =\sum_{x\in\{y,a,w\}}\sum_{s=1}^{\tau}R_{22,x(s),g^*,N}(\theta,\theta_0),
\]
\[
\begin{aligned}
R_{22,x(s),g^*,N}(\theta,\theta_0)
&=\frac1N\sum_{i=1}^N(P_{0,o(i)}-P_0^N)\bar D_{g^*,x(s)}(P)\\
&=\frac1N\sum_{i=1}^Nf^0_{x(s),g^*}(\theta,\theta_0)(C_x(i)),
\end{aligned}
\]
\[
\begin{aligned}
&f^0_{x(s),g^*}(\theta,\theta_0)(C_x(i))\\
&=\{f_{x(s),g^*}(\theta,\theta_0)(C_x(i))-P_0^Nf_{x(s),g^*}(\theta,\theta_0)\},\\
&f_{x(s),g^*}(\theta,\theta_0)(C_x(i))\\
&\equiv\frac{h^*_{c_x(s)}(C_x(i))}{\bar h_{c_x}(C_x(i))}
 \int_x\Phi_{x(s)}(x,C_x(i))
 \,\mathrm d(P^*_{x(s)}-P^*_{0,x(s)})(x\mid C_x(i)).
\end{aligned}
\]
因此，$R_{22,x(s),g^*,N}(\theta,\theta_0)$ 是 $N$ 个均值为零的随机变量
$f_i=f^0_{x(s),g^*}(\theta,\theta_0)(C_x(i))$ 的平均
$X_N=N^{-1}\sum_{i=1}^Nf_i$。这些变量是 $C_x(i)$ 的函数，每一个都在 $P\to P_0$ 时趋于零。

最后，回顾 $\Psi_J(P)=J^{-1}\sum_{j=1}^J\Psi_{g_j^*}(P)$ 和
$D^N(P)=J^{-1}\sum_{j=1}^JD^N_{g_j^*}(P)$，有
\[
\begin{aligned}
-\frac1N\sum_{i=1}^NP_{0,o(i)}\bar D(P)
={}&\Psi_J(P)-\Psi_J(P_0)\\
&+R_{21}(\theta,\theta_0)+R_{22,N}(\theta,\theta_0),
\end{aligned}
\]
其中
$R_{21}(\theta,\theta_0)=J^{-1}\sum_{j=1}^JR_{21,g_j^*}(\theta,\theta_0)$，
$R_{22,N}(\theta,\theta_0)=J^{-1}\sum_{j=1}^JR_{22,g_j^*,N}(\theta,\theta_0)$。
\end{bookresult}
